\section{计算与复现补充}
\label{sec:appendix}

\subsection{局部工作坐标的计算}
为使距离计算统一，实验固定椭球半长轴$a=6\,378\,137$米、反扁率$1/f=298.257223563$，令高度$h=0$。固定局部原点为$(\lambda_0,\phi_0)=(121.343555^\circ,31.3561015^\circ)$，没有根据候选得分移动原点。

将经纬度转成弧度，设$e^2=f(2-f)$、$N_\phi=a/\sqrt{1-e^2\sin^2\phi}$，先计算地心直角坐标：
\begin{align*}
X&=(N_\phi+h)\cos\phi\cos\lambda,\\
Y&=(N_\phi+h)\cos\phi\sin\lambda,\\
Z&=[N_\phi(1-e^2)+h]\sin\phi.
\end{align*}
减去原点的地心坐标后，取局部东、北方向：
\begin{align*}
E&=-\sin\lambda_0\Delta X+\cos\lambda_0\Delta Y,\\
N&=-\sin\phi_0\cos\lambda_0\Delta X-\sin\phi_0\sin\lambda_0\Delta Y+\cos\phi_0\Delta Z.
\end{align*}
该数学转换与PROJ文档中的地理坐标到地心坐标、固定原点地心到局部坐标的定义对应\cite{proj-cart,proj-topo}。原始坐标保持不变，图中的“相对位置”仅作平移。时空图第三轴为相对时间，不是高度，也不加入二维距离或DP误差计算。

初轮307条、30,992点的独立PROJ路径核对，最大坐标差约$4.82\times10^{-10}$工作米。该范围内1,594,527个记录内点对的ENU距离与相应椭球测地距离最大绝对差约3.503584米、最大相对差约0.015054\%。前者检查计算实现，后者描述模型近似，两者都不是仪器测量误差。实际处理阶段的敏感性检查在已检验范围内未发现登记动作变化，不能据此推广为源坐标精度已知。

\subsection{读数与实验的复算入口}
作业提供两份完成版Notebook：基础本对应轨迹预处理、参数和顺序实验，系统本对应LLM辅助与评价。两者调用共同的处理与指标实现，使用项目相对路径。默认完整复算从原始输入和固定设置重建轨迹、读数与图表，不新增模型调用；涉及模型建议时读取当时保存的真实响应，再运行确定性工具。

重新发起模型实验使用单独的显式入口，产生新的响应和结果，不能要求随机模型逐字复现旧回答。原始数据、各阶段点索引、参数表、真实建议和复现说明随工程保留；本报告用表图解释主要发现，不将文件校验和运行日志重复列入正文。
